\section{Limitations of Current Approaches}

% ---------- %

Despite notable advances in both physics-based and machine learning (ML) approaches to materials modelling, several
persistent limitations constrain the predictive fidelity, scalability, and generalisability of current methods,
particularly when applied to realistic interface systems.

This section outlines three critical domains in which these limitations manifest, motivating the hybrid methodologies
explored later in this report:

% ---------- %

\begin{itemize}
  \item \textbf{What are the current bottlenecks in interface prediction?} \\
  Interface prediction remains an inherently multiscale problem, combining combinatorial complexity at the atomic level
  with long-range elastic and electrostatic effects.

  First-principles approaches such as density functional theory (DFT) offer high accuracy but scale cubically with
  system size, making them infeasible for large or disordered systems typical of interfaces.

  Even with machine-learned interatomic potentials (MLPs), exploring configurational space, e.g. over terminations,
  interlayer slips, and misfit reconstructions, requires extensive sampling and often yields flat energy landscapes with
  many metastable minima.

  \item \textbf{Where do both ML and physics-based methods fall short?} \\
  DFT underestimates band gaps and struggles to capture emergent interfacial phenomena such as dipole formation, phase
  reconstruction, or charge transfer without recourse to hybrid functionals or many-body perturbation theory.

  Conversely, ML models, especially those trained on bulk or idealised configurations, exhibit limited transferability
  to chemically complex, strained, or disordered interfaces.

  Without embedded physical constraints or domain-specific descriptors, ML models often fail to resolve interfacial
  states, misfit dislocations, or partial intermixing.

  \item \textbf{Are there data or generalisation issues?} \\
  Interface-specific datasets are significantly less developed than those for bulk crystals, due in part to the vast
  configurational diversity of junctions and the lack of standardised repositories.

  This data sparsity hampers the ability of ML models to generalise across terminations, chemistries, or stacking
  motifs.

  Additionally, many early ML workflows lack symmetry-invariant and scale-independent representations, increasing the
  risk of overfitting to specific atom counts or lattice types.

\end{itemize}

% ---------- %

These limitations underscore the need for hybrid workflows that integrate ML potentials (e.g., MACE, ChgNet) for
large-scale screening, DFT for high-fidelity validation, and structure-generation tools such as ARTEMIS or RAFFLE for
efficient sampling of the configurational landscape.

Such frameworks aim to balance accuracy, efficiency, and transferability in modelling the behaviour of realistic
interfacial systems.

% ---------- %

\subsection{What are the current bottlenecks in interface prediction?}

% ---------- %

Despite considerable advances in computational materials science, reliable prediction of interface structures remains a
fundamental bottleneck.

This difficulty stems from the intrinsic complexity of real interfaces, which are often chemically heterogeneous,
strain-mediated, and structurally disordered.

The configurational degrees of freedom, ranging from slip vectors, surface terminations, and intermixing patterns to
stacking sequences and misfit accommodation, create a vast structural landscape that resists exhaustive exploration via
conventional methods.

% ---------- %

\paragraph{Configurational Complexity.} At the atomic scale, even between lattice-matched materials, variations in
terminations, registry shifts, and interlayer relaxations produce an exponential number of possible configurations.

The search space rapidly becomes unmanageable, especially in systems requiring large supercells to capture long-range
strain or disorder.

This renders brute-force density functional theory (DFT) calculations infeasible for comprehensive screening.

% ---------- %

\paragraph{Computational Cost.} DFT, while accurately and chemically grounded, exhibits cubic scaling with respect to
the number of electrons or basis functions, making simulations of interfaces with hundreds to thousands of atoms
prohibitively expensive.

Moreover, structural relaxations in such systems often converge poorly due to shallow energy gradients and multiple
metastable minima.

These numerical instabilities pose a challenge even for single configuration evaluations.

% ---------- %

\paragraph{Lack of Generalisable Descriptors.} Unlike bulk systems, for which symmetry and stoichiometry often impose
constraints, interfaces exhibit broken symmetries, compositional gradients, and non-uniform electrostatics.

There is no broadly applicable analogue to Anderson’s rule for predicting interfacial band alignment, necessitating
ad-hoc corrections such as $\Delta E_\Gamma$ and $\Delta E_{\mathrm{IF}}$ that work only for narrow material classes.

% ---------- %

\paragraph{Limitations of Current Structure Generation Tools.} Automated interface builders like ARTEMIS and RAFFLE
enable systematic generation of interface candidates but are typically restricted to heuristic approaches such as rigid
lattice matching or void-filling.

These tools often fail to account for elastic mismatch, surface reconstructions, or stacking instabilities.

Furthermore, the lack of integration with energetic feedback loops results in large pools of candidates requiring
expensive post-hoc evaluation.

% ---------- %

\paragraph{Transferability of Machine-Learned Potentials.} Machine-learned interatomic potentials (MLPs), such as MACE
and ChgNet, offer pathways to accelerate structural screening at near-DFT accuracy.

However, their performance is contingent upon the representativeness of the training data.

Interfacial environments frequently exhibit local chemistries and coordination motifs absent from bulk-derived datasets.

As such, model extrapolation can lead to large errors, particularly in strained, defected, or reactive regimes.

% ---------- %

\paragraph{Summary of Bottlenecks.}

\begin{itemize}

  \item The combinatorial complexity of plausible interface configurations.

  \item High computational cost and poor scaling of DFT-based relaxations.

  \item Absence of generalised physical descriptors for predicting interfacial properties.

  \item Limited adaptability of current structure generation frameworks.

  \item Transferability limitations of MLPs in interfacial chemical environments.

\end{itemize}

% ---------- %

Overcoming these bottlenecks likely requires hybrid workflows that combine MLP-based structural exploration with
targeted DFT refinements, supported by uncertainty quantification and descriptor-led screening.

Such frameworks are increasingly positioned as a pathway to scalable, realistic, and physics-informed interface
prediction.

% ---------- %

\subsection{Where do both ML and physics-based methods fall short?}

% ---------- %

Despite significant progress in both first-principles and machine learning (ML) approaches, key limitations persist in
the predictive modelling of material interfaces.

These limitations arise not only from computational constraints but also from intrinsic mismatches between real
interfacial complexity and the assumptions underpinning current models.

% ---------- %

Physics-based methods, most notably Density Functional Theory (DFT), provide high accuracy for local bonding
environments and band structure prediction. However, the computational scaling of standard Kohn-Sham DFT, typically
$\mathcal{O}(N^3)$ with respect to system size, renders large or disordered interface simulations intractable without
approximation.

Even when feasible, DFT struggles to capture extended features such as long-range strain relaxation, misfit
dislocations, and spatially diffuse defect states that dominate realistic interface behaviour.

% ---------- %

Machine learning potentials (MLPs), such as those based on graph neural networks or equivariant message passing
architectures, have emerged as promising tools to extend predictive power to larger systems at near-DFT accuracy.

Nonetheless, they too face structural limitations.

MLPs require extensive, representative training datasets, typically derived from DFT calculations, and their performance
deteriorates when applied to configurations involving unseen chemistries, surface reconstructions, or high strain
regimes.

Furthermore, many ML models lack rigorous uncertainty quantification, making it difficult to detect failure modes in
out-of-distribution predictions.

% ---------- %

Both DFT and ML approaches also fall short in capturing emergent interfacial phenomena such as charge transfer,
interface-induced dipoles, or the formation of new interfacial phases.

These effects are often strongly dependent on the local atomic registry and chemical heterogeneity, and cannot be
inferred from the properties of the isolated constituents alone.

Notably, simple models like Anderson's rule have been shown to systematically fail in 2D and 2D|3D heterostructures,
necessitating corrections such as $\Delta E_\Gamma$ and $\Delta E_\mathrm{IF}$ that explicitly depend on the
atomic-scale interface structure and electrostatics.

% ---------- %

While ML and physics-based models have complementary strengths, neither is sufficient in isolation for the reliable
prediction of interfacial behaviour.

Their limitations motivate the development of hybrid workflows, such as embedding physical constraints into ML
architectures, or coupling MLP-based screening with selective DFT refinement.

These approaches aim to achieve a balance between scalability, accuracy, and generalisability, a balance that remains at
the forefront of interface modelling challenges.

% ---------- %

\subsection{Are there data or generalisation issues?}

% ---------- %

\subsubsection{Data scarcity and imbalance}

% ---------- %

In contrast to bulk crystals, where structured repositories like the Materials Project or AFLOW provide extensive
coverage, interface-specific datasets remain sparse.

This is due to the combinatorial explosion of viable terminations, slip vectors, and orientations across material pairs,
especially in mixed-dimensional systems.

Existing databases tend to overrepresent a narrow set of lattice-matched, technologically prominent interfaces,
introducing bias and hindering the model's capacity to generalise across broader chemical or structural classes.

% ---------- %

Another critical challenge is the lack of standardisation in how interface data is represented.

Unlike bulk materials, which can be described relatively easily by a set of structural features such as lattice
constants and atomic compositions, interfaces require more complex descriptors, including local bond environments,
atomic coordination numbers, and structural motifs that change drastically depending on orientation, termination, and
strain.

As a result, the design of robust and transferable machine learning models for interfaces is further complicated by the
difficulty in developing universally applicable features that capture the intricate nature of interfacial systems.

This problem is compounded by the lack of large, consistent datasets that integrate structural, energetic, and
electronic data, making it harder for models to learn the subtle interactions that govern interface stability and
properties.

% ---------- %

Finally, the sheer diversity of possible interface configurations, especially for systems involving two-dimensional (2D)
materials, complex oxides, or disordered solids, makes high-throughput DFT calculations or other first-principles
methods impractical for exhaustive data collection.

To address these issues, there is a pressing need for hybrid data collection strategies that combine high-fidelity,
physics-based methods with more scalable, data-driven approaches, such as machine learning potential models that can
rapidly generate large datasets of interfacial configurations.

Additionally, targeted experimental datasets that validate theoretical predictions are essential for creating reliable
and comprehensive interfaces databases.

% ---------- %

\subsubsection{Transferability and chemical diversity}

% ---------- %

Machine-learned potentials (MLPs) such as MACE and ChgNet often fail to extrapolate to configurations that deviate from
the training distribution; particularly at strained, chemically heterogeneous, or electronically reconstructed
interfaces.

The emergence of site-specific hybridisation, localised charge transfer, and dipole formation presents physics not
typically encountered in bulk datasets.

Without explicit exposure to such environments during training, even state-of-the-art equivariant models exhibit
elevated uncertainty and reduced fidelity in these regimes.

% ---------- %

\subsubsection{Descriptor and representation limitations}

% ---------- %

Generalisation error is further amplified by suboptimal representation of interface-specific physics.

Many descriptor schemes, such as SOAP vectors or radial symmetry functions, are optimised for bulk periodicity and fail
to capture interfacial asymmetry, vacuum discontinuities, or broken symmetry.

Representations that do not enforce scale-invariance or rotational equivariance can also hinder model transferability to
supercells or faceted configurations.

Physically meaningful descriptors, such as interfacial dipoles, local charge density, or cleavage energy remain
underused, though they show promise for enhancing predictive robustness.

% ---------- %

\subsubsection{Label quality and training signal noise}

% ---------- %

Finally, inconsistencies in training labels present an often-overlooked source of predictive error.

Energies obtained from high-throughput DFT pipelines can vary depending on pseudopotential choice, relaxation
tolerances, or convergence thresholds.

In interfacial contexts, these uncertainties are magnified by metastability and complex relaxation pathways.

Some datasets rely on static, unrelaxed geometries or surrogate models, further undermining the reliability of ground
truth energetics.

% ---------- %

While ML offers substantial speed-ups over DFT, its utility in interface prediction is presently constrained by data
bottlenecks, representational mismatch, and transferability failures.

Addressing these limitations will require a multi-pronged strategy involving broader and more diverse training datasets,
improved descriptors attuned to interfacial physics, and integration of uncertainty quantification into ML workflows.

Hybrid schemes, where MLPs perform large-scale screening followed by selective DFT refinement, may offer a viable path
forward.